\documentclass[11pt,a4paper]{article}

\usepackage[T1]{fontenc}
\usepackage[utf8]{inputenc}
\usepackage[brazilian]{babel}
\usepackage{lmodern}
\usepackage{microtype}
\usepackage[margin=2.5cm]{geometry}
\usepackage{amsmath}
\usepackage{booktabs}
\usepackage{graphicx}
\usepackage{adjustbox}
\usepackage{placeins}
\usepackage{float}
\usepackage[hidelinks]{hyperref}

\graphicspath{{../results/figures/}{assets/}}
\setlength{\parskip}{0.5em}
\setlength{\parindent}{0pt}

% \tabela{arquivo em tables/}{legenda}
\newcommand{\tabela}[2]{%
  \begin{table}[H]
    \centering
    \caption{#2}
    \label{tab:#1}
    \begin{adjustbox}{max width=\textwidth}
      \input{tables/#1}
    \end{adjustbox}
  \end{table}%
}

% \figura{arquivo em results/figures/}{legenda}{largura relativa ao texto}
\newcommand{\figura}[3]{%
  \begin{figure}[H]
    \centering
    \includegraphics[width=#3\textwidth]{#1}
    \caption{#2}
    \label{fig:#1}
  \end{figure}%
}

\title{\includegraphics[width=2.6cm]{brasao_uece}\\[1em]
  Experimento dirigido: MLP com backpropagation na base Dry Bean\\
  \large Avaliação de Inteligência Computacional, 2026.2}
\author{[Samuel William Silva Almeida\\
  \small Matrícula: 1631353]}
\date{Novembro de 2026}

\begin{document}

\maketitle

\begin{abstract}
  Este relatório investiga uma rede MLP treinada por backpropagation na classificação de sete variedades de
  feijão a partir de 16 atributos geométricos. Foram avaliados inicialização, taxa de aprendizado, momento,
  topologia, quantidade e qualidade dos dados, atributos e, por fim, a generalização por validação cruzada.
  A configuração final, com camadas ocultas 80$\times$80, LR = 0,01, momento = 0,7 e 10 atributos de forma, obteve
  Medida-F macro de $0{,}9436 \pm 0{,}0085$ em validação cruzada com 10 partições. A variação entre partições
  superou a variação entre as 30 topologias testadas, e a regressão logística ficou a 0,0070 da MLP.
\end{abstract}

\clearpage
\tableofcontents
\clearpage

\section{Introdução e metodologia comum}

A tarefa é classificar grãos de feijão seco em sete variedades (Barbunya, Bombay, Cali, Dermason, Horoz, Seker e
Sira) a partir de 16 atributos numéricos extraídos de imagens \cite{koklu2020,uci}. Todo o experimento está em
\texttt{notebooks/mlp\_dry\_bean.ipynb}, e as tabelas e figuras deste texto são geradas por ele.

O enunciado traz três sobras de versões anteriores: cita o Fashion-MNIST na Questão 1, chama de ``Questão 5'' e
``Questão 6'' os relatos das seções 6 e 7, e chama de ``Questão 4'' a validação inicial. Este relatório segue os títulos
das seções e trata a base como Dry Bean em todas as etapas.

\subsection{Dados}

A base original tem 13.611 exemplos, sem valores ausentes. Havia 68 linhas duplicadas, removidas antes da divisão para
que a mesma linha não caísse no treino e no teste ao mesmo tempo (Tabela~\ref{tab:dados_visao_geral}). As classes são
desbalanceadas: Dermason tem 3.546 exemplos e Bombay, 522 (Figura~\ref{fig:dados_classes}). Os atributos são muito
correlacionados, com $|r|$ de 0,999 entre \texttt{Compactness} e \texttt{ShapeFactor3}, e de aproximadamente 1,00 entre
\texttt{Area} e \texttt{ConvexArea} (Figura~\ref{fig:dados_correlacao}).

\tabela{dados_visao_geral}{Exemplos em cada etapa da preparação. A divisão é estratificada e fixa (70/15/15).}
\figura{dados_classes}{Exemplos por classe, após a remoção de duplicatas.}{0.65}
\figura{dados_correlacao}{Correlação de Pearson entre os 16 atributos.}{0.8}

\subsection{Modelo e pré-processamento}

O modelo é o \texttt{MLPClassifier} do scikit-learn \cite{sklearn}, com SGD, momento clássico (sem Nesterov), lotes de 200
exemplos e regularização L2 com $\alpha = 10^{-4}$, salvo onde indicado. As camadas ocultas usam ReLU; a de saída usa
softmax e a perda é a entropia cruzada categórica. A classe aplica softmax e entropia cruzada automaticamente em
problemas multiclasse e converte os rótulos internamente em one-hot, cumprindo a codificação recomendada no enunciado.
O z-score é ajustado apenas no conjunto de treino de cada execução e aplicado depois aos demais, o que evita vazamento.

\subsection{Laço de treinamento}

Para registrar métricas a cada época, o treino usa \texttt{partial\_fit}. A alternativa sugerida na documentação,
\texttt{warm\_start=True} com \texttt{max\_iter=1}, recria o otimizador a cada chamada e zera a velocidade do momento no início
de cada época. Além disso, uma seed inteira repete o embaralhamento em todas as épocas, por isso a seed é passada como
um objeto \texttt{RandomState}.

A Tabela~\ref{tab:q2_laco_treino} e a Figura~\ref{fig:q2_laco_treino} medem esses efeitos. Depois de 200 épocas, a perda
com \texttt{warm\_start} e seed inteira fica 3\% acima da obtida com \texttt{partial\_fit} em LR = 0,01 e 25\% acima em
LR = 0,1. O efeito é moderado porque cada época tem 48 lotes e, com momento 0,9, a velocidade se recompõe em cerca de 10.

\tabela{q2_laco_treino}{Perda de treino na época 200 para três laços de treinamento (três seeds).}
\figura{q2_laco_treino}{Curvas de perda dos três laços, a partir da quinta época.}{0.85}

\subsection{Critérios e métricas}

\begin{itemize}
  \item \textbf{Parada:} perda de treino $\le 0{,}001$ ou limite de épocas. Como esse valor raramente é atingido, a velocidade
    também é medida pelas épocas até a perda de treino chegar a 0,20.
  \item \textbf{Perda:} entropia cruzada calculada no conjunto inteiro ao fim de cada época, sem o termo L2.
  \item \textbf{Medida-F, precisão e revocação:} médias macro, em que cada classe pesa igual, inclusive a minoritária.
  \item \textbf{Repetições:} três ou cinco seeds por configuração. Os resultados são média e desvio-padrão (dp).
  \item \textbf{Tempo:} os treinos rodam em paralelo, e os tempos medidos nessa condição não são comparáveis. Uma medição
    à parte roda um treino por vez, em uma thread, por 30 épocas, e registra o mínimo de cinco repetições intercaladas. Os
    valores são de uma única máquina e servem para comparar configurações entre si.
  \item \textbf{Uso do teste:} até a Questão 5 toda escolha usa a validação. As Questões 4 e 5 reportam o teste porque o
    enunciado pede, sem decidir nada com ele.
\end{itemize}

\FloatBarrier

\section{Questão 1: pesos iniciais e funções de ativação}

\subsection*{a. Descrição da configuração experimental}

A rede tem 16 entradas, duas camadas ocultas com 50 neurônios cada e 7 saídas. Esse tamanho está dentro da faixa
explorada na Questão 3 e serve de ponto de partida. As camadas ocultas usam ReLU, que não satura para ativações positivas
e costuma convergir bem com a inicialização usada pela classe. A saída usa softmax, que produz uma distribuição de
probabilidade sobre as classes. A perda é a entropia cruzada categórica, a escolha padrão para esse tipo de saída.

O otimizador é SGD com momento = 0,9 e LR = 0,01, por até 300 épocas. Usar SGD desde o início mantém esta etapa comparável
com a Questão 2. Foram feitos cinco treinamentos com seeds diferentes e, para justificar a ReLU, os mesmos cinco com tanh e
logistic. Só o conjunto de treino é usado.

\subsection*{b. Resultados}

\tabela{q1_seeds}{Cinco treinamentos com ReLU, no conjunto de treino.}
\tabela{q1_ativacoes}{Comparação das funções de ativação (média de cinco seeds).}
\figura{q1_curvas_seeds}{Perda e acurácia de treino com ReLU, por seed, a partir da quinta época.}{0.8}
\figura{q1_curvas_ativacoes}{Perda (em escala logarítmica) e acurácia de treino por função de ativação.}{0.8}

\subsection*{c. Análise e discussão}

Os cinco treinamentos com ReLU são quase coincidentes: a perda final varia de 0,1460 a 0,1502, e a acurácia de 0,9441 a
0,9463. A inicialização afeta só as primeiras épocas, e o desvio final é pequeno (0,0015 na perda, 0,0009 na acurácia).
As curvas seguem caindo devagar na época 300, com oscilações pequenas, sem sinal de instabilidade.

Quanto à ativação, ReLU e tanh chegam a resultados próximos (perda 0,1480 e 0,1520). A logistic demora muito mais: são
239 épocas até a perda de 0,20, contra 12 da ReLU, e fica com a pior perda final (0,1947). Nas primeiras épocas, sua
acurácia fica parada em cerca de 0,26, a proporção da classe majoritária.

Sobre over e underfitting, os dados de treino sozinhos só permitem uma conclusão parcial. Não há underfitting evidente,
pois a acurácia passa de 0,94 em todas as seeds. O overfitting não pode ser avaliado sem um conjunto separado, o que é
feito a partir da próxima questão.

\subsection*{d. Conclusão da etapa}

O modelo aprende os padrões básicos, com acurácia de treino de 0,945. O treinamento é estável em relação aos pesos iniciais,
e a ReLU é adequada: converge tão bem quanto a tanh, é muito mais rápida que a logistic e não apresentou nenhuma
instabilidade nessa configuração. A topologia 50$\times$50 com ReLU segue como base da Questão 2.

\FloatBarrier

\section{Questão 2: taxa de aprendizado e termo momento}

\subsection*{a. Descrição do experimento}

Foi feita uma grade de 12 combinações: LR em \{0,001; 0,01; 0,1; 0,5\} e momento em \{0,5; 0,7; 0,9\}, com três seeds cada.
A rede é a da Questão 1 (50$\times$50, ReLU, softmax, entropia cruzada), treinada por SGD com momento clássico. O critério de
parada é perda de treino $\le 0{,}001$, com limite de 500 épocas.

O treino atualiza os pesos, e a validação apenas acompanha a generalização. A regra de escolha foi fixada antes de rodar a grade:
entre as combinações sem divergência e com perda final de validação até 5\% acima da melhor, fica a mais rápida até a perda de
treino de 0,20. A tolerância existe porque diferenças tão pequenas ficam perto da variação entre seeds.

\subsection*{b. Resultados}

\tabela{q2_grade}{Grade de LR e momento (média de três seeds, em 500 épocas). ``Atingiu'' e ``Divergiu'' contam execuções.}
\figura{q2_heatmap_perdas}{Perda final de treino e de validação. A cor usa escala logarítmica.}{0.75}
\figura{q2_heatmap_velocidade}{Épocas até a perda de 0,20 e desvio da perda entre seeds. A célula vazia é a combinação que
nunca chegou a 0,20.}{0.75}
\figura{q2_curvas}{Curvas de perda em escala logarítmica. A linha tracejada marca o critério de 0,001.}{1}

A combinação escolhida foi \textbf{LR = 0,01 e momento = 0,7}: perda de treino 0,1584, perda de validação 0,1908 e 28 épocas
até a perda de 0,20. As candidatas dentro da tolerância foram LR = 0,001 com momento 0,9 (91 épocas), LR = 0,01 com momento 0,5
(45 épocas) e a própria escolhida.

\subsection*{c. Análise e discussão}

\textbf{Critério de 0,001.} Nenhuma combinação atingiu esse valor em 500 épocas. Por isso, a coluna ``Épocas até 0,001''
fica vazia, e a velocidade foi comparada pelo limiar de 0,20.

\textbf{Taxa de aprendizado muito baixa.} Com LR = 0,001 e momento 0,5, são 447 épocas até a perda de 0,20, e a perda de
validação ainda está caindo no fim do treino.

\textbf{Taxa de aprendizado alta.} Com LR $\ge$ 0,1, a perda de treino cai mais e fica ruidosa, mas a de validação sobe:
LR = 0,1 com momento 0,9 termina em 0,0735 no treino e 0,5709 na validação. As menores perdas de validação ao longo do treino são
parecidas em quase toda a grade (de 0,188 a 0,196, à exceção do colapso descrito abaixo e de LR = 0,001 com momentos 0,5 e 0,7), mas são atingidas cedo
quando a LR é alta. Nesse caso o problema é overfitting rápido, e a capacidade de aprender é a mesma.

\textbf{Colapso.} Com LR = 0,5 e momento 0,9, a perda fica em 1,84 nas três seeds e a Medida-F é 0,0593. A entropia das classes
é 1,8338, e prever sempre a classe majoritária dá Medida-F de 0,0593, ou seja, a rede aprendeu apenas a proporção das
classes. Como os pesos continuam finitos, nenhuma execução foi contada como divergência numérica. A causa provável é a morte de
neurônios ReLU após os primeiros passos grandes, que o experimento não verificou diretamente.

\textbf{Efeito do momento.} O número de épocas até a perda de 0,20 acompanha a taxa efetiva $\eta/(1-m)$. O produto das épocas por essa
taxa fica entre 0,89 e 0,93 nas cinco combinações com taxa efetiva de até 0,033. Assim, LR = 0,001 com momento 0,9 (taxa efetiva 0,01)
leva 91 épocas, e LR = 0,01 com momento 0,5 (taxa efetiva 0,02) leva 45. Nessa faixa, o momento acelera o treino quase como aumentar a
LR. A estabilidade, por outro lado, piora com os dois: o maior desvio da perda entre seeds ocorre em
LR = 0,1 com momento 0,9 (0,0265).

\subsection*{d. Conclusão da etapa}

LR = 0,01 com momento = 0,7 teve a menor perda final de validação da grade e, entre as candidatas, a convergência mais rápida,
sem oscilação visível. Essa combinação fica fixa na Questão 3. Momento e taxa de aprendizado aceleram o treino de forma parecida;
acima de uma taxa efetiva de aproximadamente 0,1, o ganho de velocidade vem acompanhado de overfitting rápido e, no extremo, de colapso.

\FloatBarrier

\section{Questão 3: número de camadas e de neurônios}

\subsection*{a. Descrição da configuração experimental}

Foram testadas 30 topologias: 1, 2 e 3 camadas ocultas iguais, de 10 a 100 neurônios por camada, em passos de 10, com três seeds cada.
Ficam fixos ReLU, softmax, entropia cruzada, SGD com LR = 0,01 e momento = 0,7, e $\alpha = 10^{-4}$. A parada é perda de treino
$\le 0{,}001$ ou 500 épocas. O teste não é usado. O tempo total da tabela é o tempo por época multiplicado pelas épocas executadas.

\subsection*{b. Resultados}

\tabela{q3_topologias}{As 30 topologias (média de três seeds, em 500 épocas). ``Diferença'' é a perda de validação menos a de treino.}
\figura{q3_heatmaps}{Medida-F, perdas e tempo por época para cada topologia.}{0.72}
\figura{q3_curvas}{Perda de treino e de validação das quatro melhores redes e de dois extremos (10 e 100$\times$100$\times$100), a partir da
época 20.}{0.85}

Com LR = 0,01, nenhuma topologia atinge a perda de 0,001. Para examinar esse critério, um caso de contraste à parte da grade usa
a maior rede com LR = 0,1 e momento = 0,9 (Tabela~\ref{tab:q3_contraste} e Figura~\ref{fig:q3_contraste}).

\tabela{q3_contraste}{Contraste: 100$\times$100$\times$100 com LR = 0,1 e momento = 0,9. Fora da grade e da seleção.}
\figura{q3_contraste}{Perda de treino e de validação no caso de contraste.}{0.6}

\subsection*{c. Análise e discussão}

\textbf{A topologia muda pouco o desempenho.} A Medida-F de validação vai de 0,9322 (10 neurônios) a 0,9384 (40$\times$40$\times$40), uma faixa de 0,0062,
da mesma ordem dos desvios entre seeds (até 0,0038). Uma rede com uma camada de 70 neurônios (0,9371) fica muito perto de
80$\times$80 (0,9380) e leva 21,1~ms por época, contra 34,2~ms.

\textbf{A complexidade reduz a perda de treino, mas não a de validação.} De 10 neurônios a 100$\times$100$\times$100, a perda de treino cai de 0,1835
para 0,1359. A de validação melhora só até redes médias e fica em torno de 0,19. A diferença entre as duas é de 0,02 com uma camada,
0,03 com duas e chega a 0,06 com três.

\textbf{Underfitting.} A rede de 10 neurônios tem as maiores perdas (0,1835 e 0,2088), o que é um underfitting leve, já que ela ainda alcança Medida-F
de 0,9322. Nenhuma das redes testadas falha em aprender.

\textbf{Overfitting.} Na grade, os sinais são moderados. Com 1 e 2 camadas, a validação ainda melhora por volta da época 450 a 500. As redes com
3 camadas e pelo menos 30 neurônios atingem o mínimo entre as épocas 269 e 345 e depois sobem levemente. No contraste, o
overfitting é claro: as três seeds atingem o critério de 0,001 entre as épocas 363 e 394, com acurácia de treino de pelo menos 0,9999, enquanto a
perda de validação sobe de cerca de 0,19 (nas épocas 24 a 32) para mais de 0,79. A Medida-F de validação cai bem menos, para cerca de 0,926: a rede
continua classificando razoavelmente, com probabilidades excessivamente confiantes. O critério de 0,001 mede, portanto, a capacidade de memorizar o treino,
e atingi-lo não indica boa generalização nesta base.

\textbf{Tempo.} O tempo por época cresce com o número de pesos: 15,8~ms com 10 neurônios, 41,0~ms com 100$\times$100 e 57,8~ms com
100$\times$100$\times$100. Em 500 épocas, o total vai de 7,9~s a 28,9~s. O ponto de equilíbrio entre generalização e custo está nas redes de uma camada com
60 a 90 neurônios ou nas de duas camadas com 60 a 80. Os ganhos das redes maiores ficam dentro da variação entre seeds.

\subsection*{d. Seleção das melhores topologias}

A regra, fixada antes de rodar, ordena as redes pela média menos o desvio da Medida-F de validação, o que pontua desempenho e estabilidade.

\tabela{q3_top4}{As quatro topologias selecionadas.}

\begin{itemize}
  \item \textbf{80$\times$80:} a melhor pontuação, com Medida-F de 0,9380 e desvio de só 0,0006.
  \item \textbf{90$\times$90:} a menor perda final de validação entre as selecionadas (0,1853), com desvio de 0,0008.
  \item \textbf{80$\times$80$\times$80:} a representante com três camadas, com Medida-F de 0,9376, embora com diferença de perda maior (0,0545) e o maior custo do grupo (45,8~ms por época).
  \item \textbf{60$\times$60$\times$60:} o menor desvio da grade (0,0002), com custo menor que o da rede anterior (36,9~ms).
\end{itemize}

A rede 40$\times$40$\times$40 tem a maior média (0,9384), mas fica de fora pelo desvio de 0,0028. A ordem das quatro é pouco sensível em termos práticos, dada a proximidade dos valores.

\subsection*{e. Conclusão da etapa}

Nesta base, com LR = 0,01, o número de camadas e de neurônios tem efeito pequeno na generalização. O aumento de capacidade reduz a perda de treino e
aumenta o custo, e os sinais de overfitting crescem com a profundidade. As quatro redes seguem para a Questão 6, e a 80$\times$80 é a base das Questões 4 e 5.

\FloatBarrier

\section{Questão 4: quantidade e qualidade dos dados de treino}

\subsection*{a. Descrição do experimento}

Rede 80$\times$80 com LR = 0,01, momento = 0,7, ReLU, softmax e entropia cruzada, por até 500 épocas e com cinco seeds. Validação e teste permanecem fixos em todas as
condições, de modo que os números são comparáveis. A seed define tanto a amostra quanto os pesos iniciais.

\begin{itemize}
  \item \textbf{Quantidade:} 5, 10, 20, 40, 60, 80 e 100\% do conjunto de treino, com amostragem estratificada. As frações de 5 e 10\% foram acrescentadas às sugeridas para localizar melhor a região de queda.
  \item \textbf{Qualidade dos rótulos:} 5, 10 e 20\% dos rótulos de treino trocados por outra classe sorteada. Validação e teste mantêm os rótulos corretos.
  \item \textbf{Representatividade:} amostragem estratificada contra aleatória, com 5 e 20\% dos dados.
\end{itemize}

\subsection*{b. Resultados}

\tabela{q4_quantidade}{Efeito da quantidade de exemplos de treino (média de cinco seeds).}
\figura{q4_generalizacao}{Curvas de generalização: perda, acurácia e Medida-F contra o número de exemplos de treino. As barras indicam $\pm$ 1 dp.}{1}
\figura{q4_tempo}{Tempo por época contra o número de exemplos de treino.}{0.5}
\figura{q4_curvas_fracoes}{Curvas de perda para cada fração de dados, a partir da época 50.}{0.8}
\tabela{q4_ruido}{Efeito de rótulos errados no treino, com 100\% dos dados.}
\figura{q4_curvas_ruido}{Curvas de perda com rótulos de treino corrompidos, a partir da época 10.}{0.8}
\tabela{q4_amostragem}{Amostragem estratificada contra aleatória.}

\subsection*{c. Análise e discussão}

\textbf{Quantidade.} A generalização melhora de forma contínua, com retorno decrescente. De 5 a 40\%, a Medida-F de teste vai de 0,9176 a 0,9305. De 40 a 100\%, ganha só mais 0,0053 e
chega a 0,9358. A perda de treino pouco muda (0,14 a 0,15), enquanto a de validação cai de 0,2666 para 0,1894: a diferença entre as duas passa de 0,122 para 0,035. A perda de
validação não satura dentro da faixa testada, embora a Medida-F quase sature. Isso indica que mais dados ainda melhoram a qualidade das probabilidades previstas, mesmo com poucos erros de classe a menos.

Há uma ressalva de desenho: como o limite é de 500 épocas, frações menores recebem menos atualizações dos pesos (1.500 com 5\% e 24.000 com 100\%). Nas curvas, a perda de treino de todas
as frações converge para o mesmo nível antes da época 500, o que sugere que a falta de atualizações não domina o resultado. O efeito, porém, não foi separado.

\textbf{Qualidade dos rótulos.} O ruído afeta mais a perda do que a classificação. Com 20\% de rótulos errados, a perda de validação dobra (de 0,1894 para 0,3891), enquanto a Medida-F de
teste cai de 0,9358 para 0,9314. Como o ruído é uniforme, a classe correta continua sendo a mais frequente em cada região do espaço de atributos, e a decisão muda pouco. O que piora
é a confiança das probabilidades. Um ruído concentrado em classes parecidas, como Sira e Dermason, provavelmente causaria mais dano, mas isso não foi testado.

\textbf{Representatividade.} A estratificação quase não muda a média, mas reduz a variabilidade em amostras muito pequenas: com 5\% dos dados, o desvio da Medida-F de teste é de 0,0028 na
estratificada e de 0,0082 na aleatória. Com 20\%, os desvios são quase iguais (0,0030 e 0,0029) e os dois métodos se equivalem.

\textbf{Tempo.} O tempo por época é linear no número de exemplos: 2,6~ms com 473 e 35,5~ms com 9.479. Em 500 épocas, são cerca de 1,3~s e 17,8~s.

\subsection*{d. Conclusão da etapa}

A rede já generaliza bem com poucos dados, e o ganho a partir de cerca de 40\% do treino (3.791 exemplos) é pequeno em Medida-F. O custo cresce de forma linear e é baixo nesta base, portanto o melhor
equilíbrio é usar todos os dados disponíveis, que dão a menor perda e o menor desvio no teste. Com um limite real de coleta, 40 a 60\% já entregariam uma Medida-F de teste a cerca de 0,005 ou menos da obtida com todos os dados.
A rede mostrou boa tolerância a ruído uniforme nos rótulos, com perda de qualidade concentrada nas probabilidades.

\FloatBarrier

\section{Questão 5: quantidade, qualidade e sensibilidade dos atributos}

\subsection*{a. Descrição do experimento}

A rede é a mesma da Questão 4 (80$\times$80, LR = 0,01, momento = 0,7, ReLU, softmax, entropia cruzada, até 500 épocas), com três seeds por subconjunto.

\textbf{Critérios.} A relevância é medida pela queda da Medida-F de validação ao embaralhar o atributo (média de três redes). A redundância é $|r| > 0{,}95$ no treino: percorrendo os atributos do mais ao menos importante,
cada um só entra se não for redundante com os já mantidos. O ruído é representado por colunas gaussianas artificiais.

\textbf{Conjuntos testados.}
\begin{itemize}
  \item todos os 16 atributos;
  \item os 12, 8 e 4 mais importantes por permutação;
  \item 10 atributos não redundantes;
  \item agrupamento físico: só tamanho (6 atributos) e só forma (10 atributos);
  \item todos, mais 5 colunas de ruído (21);
  \item 8 atributos sorteados, um sorteio por seed em cinco seeds;
  \item remoção de um atributo por vez (16 conjuntos de 15).
\end{itemize}

A regra de escolha, fixada antes de rodar, é: entre os conjuntos nomeados (sem ruído e sem sorteio), o menor cuja Medida-F de validação fique a até um desvio da melhor. O teste é só reportado.

\subsection*{b. Resultados}

\tabela{q5_subconjuntos}{Resultados por conjunto de atributos. Com 8 atributos sorteados, os conjuntos mudam a cada seed, e por isso o tempo não foi medido.}
\figura{q5_subconjuntos}{Medida-F de validação por conjunto, com $\pm$ 1 dp.}{0.7}
\figura{q5_curvas}{Curvas de perda e de acurácia para seis conjuntos, a partir da época 50.}{0.8}
\tabela{q5_importancia}{Importância por permutação na validação.}
\figura{q5_importancia}{Queda da Medida-F ao embaralhar cada atributo.}{0.7}
\tabela{q5_remocao}{Remoção de um atributo por vez, em relação aos 16 atributos (Medida-F de validação de 0,9380).}
\figura{q5_remocao}{Variação da Medida-F ao remover cada atributo. As barras mostram o dp do conjunto com a remoção.}{0.7}

\subsection*{c. Análise e discussão}

\textbf{Remoção e inclusão.} Reduzir de 16 para 10 atributos não custa nada mensurável, se a escolha for boa: ``Só forma'' tem Medida-F de 0,9384 e ``Sem redundantes'', 0,9365, contra 0,9380 de todos. Cortar mais custa pouco: 0,9318 com 4 atributos,
mas a convergência fica mais lenta (112 épocas até a perda de 0,20, contra 27). Conjuntos escolhidos sem critério são piores: os 8 sorteados têm média de 0,9223 e desvio de 0,0117, o que mostra que a qualidade pesa mais que a quantidade de atributos.

\textbf{Tamanho e forma.} Os 6 atributos de tamanho sozinhos caem para 0,9197, o pior resultado, e nunca chegam à perda de treino de 0,20. Os de forma igualam o conjunto completo, dentro da variação entre seeds. Isso se explica em parte pela definição dos atributos: \texttt{ShapeFactor1} e
\texttt{ShapeFactor2} são comprimentos de eixo divididos pela área \cite{koklu2020}, e por isso o conjunto de forma também carrega informação de escala.

\textbf{Atributos redundantes.} Os atributos de tamanho são quase cópias uns dos outros. Removê-los um a um praticamente não altera a Medida-F, e as variações ficam em torno de um desvio.

\textbf{Atributos ruidosos.} As 5 colunas de ruído reduzem a perda de treino (de 0,1539 para 0,1258), mas aumentam a de validação (de 0,1882 para 0,2179) e reduzem a Medida-F de 0,9380 para 0,9325. A rede usa a capacidade extra para memorizar ruído, e a curva de validação
passa a subir depois de cerca de 150 épocas. Esse é um caso de overfitting provocado apenas pela entrada.

\textbf{Permutação contra remoção.} Os dois métodos discordam, e a discordância é informativa. \texttt{ConvexArea} perde 0,1028 quando embaralhado, mas sua remoção não piora o resultado (variação de +0,0014). O motivo é a correlação: embaralhar um atributo quase igual a \texttt{Area} cria combinações
que nunca existiram no treino, e a rede erra; removê-lo não muda a informação disponível. Logo, a permutação mede a dependência do modelo treinado, e a remoção mede se a informação é insubstituível. Por esse segundo critério, os atributos mais relevantes são
\texttt{ShapeFactor4} ($-0{,}0075$), \texttt{Extent} ($-0{,}0037$) e \texttt{Solidity} ($-0{,}0026$). \texttt{ShapeFactor4} é o único que se destaca nos dois métodos. Como cada remoção usa só três seeds, as variações menores que cerca de 0,001 devem ser lidas como empate.

\textbf{Dimensionalidade e tempo.} O tempo por época varia pouco: 32,7~ms com 6 atributos e 36,7~ms com 21, porque os pesos de entrada são uma pequena parte do total nesta rede. A dimensionalidade não afetou a estabilidade dos conjuntos bem escolhidos. A estabilidade dependeu de quais atributos entram.

\subsection*{d. Conclusão da etapa}

O conjunto com melhor equilíbrio é \textbf{``Só forma''} (10 atributos: \texttt{AspectRation}, \texttt{Eccentricity}, \texttt{Extent}, \texttt{Solidity}, \texttt{roundness}, \texttt{Compactness} e os quatro \texttt{ShapeFactor}). Ele teve a maior Medida-F de validação entre os conjuntos nomeados
e empata com os 16 atributos usando 6 a menos, com tempo por época praticamente igual. O conjunto ``Sem redundantes'' ficou 0,0001 abaixo do corte da regra e é uma alternativa equivalente. O ganho principal é a simplicidade da entrada, já que a economia de tempo é pequena.

\FloatBarrier

\section{Questão 6: validação inicial no conjunto de teste}

\subsection*{a. Descrição do experimento}

\textbf{Redes.} As quatro topologias da Questão 3 (80$\times$80, 90$\times$90, 80$\times$80$\times$80 e 60$\times$60$\times$60), com ReLU, softmax, entropia cruzada, SGD com LR = 0,01 e momento = 0,7, usando os 10 atributos de forma da Questão 5. O critério de seleção vem da Questão 3:
média menos desvio da Medida-F de validação.

\textbf{Ajustes, decididos na validação.} Testou-se $\alpha$ em $\{10^{-4}, 10^{-3}, 10^{-2}\}$ para cada topologia (três seeds), e ficou o valor de menor perda mínima de validação. O limite de épocas passou a ser a época média dessa perda mínima, para interromper o treino antes de o overfitting
avançar (Tabela~\ref{tab:q6_ajustes}).

\textbf{Avaliação.} Cada rede ajustada foi treinada com cinco seeds no conjunto de treino. O conjunto de teste entra para comparar a generalização, com a perda de entropia cruzada, a acurácia e a Medida-F.

\subsection*{b. Resultados}

\tabela{q6_ajustes}{Ajuste de regularização nos atributos de forma, na validação (três seeds).}
\tabela{q6_finalistas}{As quatro redes ajustadas, em treino e teste (média de cinco seeds). A coluna dp refere-se à Medida-F de teste.}
\figura{q6_curvas_perda}{Perda de treino e de teste das quatro redes, a partir da época 20.}{0.75}
\figura{q6_curvas_acuracia}{Acurácia de treino e de teste das quatro redes, a partir da época 20.}{0.75}
\tabela{q6_por_classe}{Rede vencedora no teste, por classe (seed de Medida-F mediana).}
\figura{q6_confusao}{Matriz de confusão da rede vencedora no teste. A cor é a taxa por classe real, e o número é a contagem.}{0.55}

Pela regra de maior média menos desvio, a melhor rede é \textbf{80$\times$80} ($\alpha = 10^{-4}$, 481 épocas), com pontuação de 0,9348. Ela tem Medida-F de teste de 0,9359, com o segundo menor desvio (0,0011).

\subsection*{c. Análise e discussão}

\textbf{As quatro redes são equivalentes na prática.} As Medidas-F de teste ficam entre 0,9345 e 0,9361, uma diferença máxima de 0,0016, menor que o desvio entre seeds de duas delas. A rede 90$\times$90 tem a maior média (0,9361), mas com o dobro do desvio de 80$\times$80. A escolha segue a regra, e os dados não sustentam a afirmação de que uma
delas é superior. A 80$\times$80 também tem tempo por época menor que o da 90$\times$90 e o da 80$\times$80$\times$80 (Questão 3).

\textbf{Treino e teste.} A rede vencedora tem Medida-F de 0,9540 no treino e 0,9359 no teste, e perda de 0,1542 e 0,2139. Essa diferença existe em todas as redes e se mantém estável ao longo das épocas, sem aumento expressivo da perda de teste. Há um overfitting leve e nenhum sinal de underfitting. O teste também parece mais difícil que a validação para esta divisão específica: a perda
é de 0,2139 no teste e cerca de 0,187 na validação, e a Questão 7 ajuda a avaliar essa diferença.

\textbf{Ajustes.} A regularização quase não teve efeito: dentro de cada topologia, a perda mínima de validação varia no máximo 0,0006 entre os valores de $\alpha$. Valores maiores de $\alpha$ foram escolhidos em duas redes, mas com diferenças que não têm peso prático. A interrupção do treino importa um pouco mais nas redes de 3 camadas, em que o mínimo surge antes
(269 e 342 épocas). O maior ajuste foi usar os atributos de forma, que mantêm o desempenho com uma entrada menor.

\textbf{Erros por classe.} Bombay é classificada sem erros, com 78 de 78 exemplos. Sira tem a menor Medida-F (0,8695), e a maior parte dos erros está entre Sira e Dermason: 43 exemplos de Dermason previstos como Sira, e 28 no sentido oposto. A confusão seguinte é de 12 exemplos de Barbunya previstos como Cali. O erro restante se deve, portanto, à sobreposição
entre variedades parecidas, o que mais capacidade de rede não parece resolver.

\subsection*{d. Conclusão da etapa}

A configuração final é a rede \textbf{80$\times$80} com ReLU, softmax, entropia cruzada, SGD com LR = 0,01, momento = 0,7, $\alpha = 10^{-4}$, 481 épocas e 10 atributos de forma. Ela combina Medida-F de teste entre as mais altas e desvio pequeno, com tempo por época menor que o da 90$\times$90, a rede de maior média. Como as finalistas são muito próximas em uma única divisão dos dados,
a validação cruzada da próxima etapa é necessária para estimar a generalização sem depender dessa divisão.

\FloatBarrier

\section{Questão 7: validação cruzada k-fold}

\subsection*{a. Descrição do experimento}

\textbf{Valor de k.} Foi usado k = 10. Cada teste tem cerca de 1.354 exemplos, dos quais cerca de 52 são da menor classe (Bombay), o que ainda dá uma estimativa razoável por classe. Cada treino usa 90\% dos dados, perto do que o modelo final usaria. O custo é baixo para esta base.

\textbf{Configuração.} Rede 80$\times$80, ReLU, softmax, entropia cruzada, SGD com LR = 0,01 e momento = 0,7, $\alpha = 10^{-4}$, 481 épocas, 10 atributos de forma.

\textbf{Procedimento.} A divisão usa \texttt{StratifiedKFold} sobre a base completa sem duplicatas (13.543 exemplos), o que mantém a proporção das classes em cada partição. Em cada rodada, o z-score é ajustado só nas nove partições de treino. A seed dos pesos é a mesma em todas as rodadas, de modo que a variação observada vem das partições.
A regressão logística, nas mesmas partições, serve de referência linear.

\subsection*{b. Resultados}

\tabela{q7_particoes}{Métricas de teste por partição e perda de treino final.}
\tabela{q7_estatisticas}{Média, variância e extremos das métricas entre as dez partições.}
\figura{q7_particoes}{Medida-F de teste por partição. A linha tracejada marca a média.}{0.65}
\figura{q7_curvas_perda}{Perda de treino e de teste na pior (1), na mediana (2) e na melhor (6) partição, a partir da época 20.}{1}
\figura{q7_curvas_acuracia}{Acurácia de treino e de teste nas mesmas partições, a partir da época 20.}{1}
\tabela{q7_por_classe}{Métricas por classe, com as previsões das dez partições reunidas.}
\figura{q7_confusao}{Matriz de confusão com as previsões das dez partições reunidas.}{0.55}
\tabela{q7_baseline}{MLP e regressão logística nas mesmas partições.}

\subsection*{c. Análise e discussão}

\textbf{Resultado médio e variabilidade.} A Medida-F é de $0{,}9436 \pm 0{,}0085$ (variância de 0,00007), a acurácia é de $0{,}9329 \pm 0{,}0101$ e a perda, de $0{,}1845 \pm 0{,}0260$. Os valores de Medida-F vão de 0,9291 a 0,9572.

\textbf{Robustez.} A rede se comporta de modo consistente: nenhuma partição fica longe das demais, e o desvio é menor que 1\% da média. Por outro lado, esse desvio é maior que toda a faixa de Medida-F entre as 30 topologias da Questão 3 (0,0062). Assim, a escolha de quais exemplos caem no teste pesa mais no resultado do que a escolha da topologia dentro da faixa estudada.
O resultado no teste da Questão 6 (0,9359) está dentro da faixa e perto do limite inferior de um desvio (0,9351), o que confirma que aquela divisão era um pouco mais difícil que a média.

\textbf{Flutuações.} A perda de treino quase não muda entre partições (0,1561 a 0,1630), enquanto a de teste vai de 0,1447 a 0,2417. Na partição 6, a perda de teste é até menor que a de treino. A variação vem, portanto, de quais exemplos caem no teste, e o treinamento em si é estável. Quanto à causa, o provável é o número de exemplos das fronteiras entre classes parecidas presentes em cada partição, mais do que
o desbalanceamento em si: a estratificação fixa as proporções, e a classe menor, Bombay, é quase sempre acertada (Medida-F de 0,9990). Essa explicação não foi testada partição a partição.

\textbf{Erros por classe.} A matriz reunida confirma a Questão 6. Sira tem a menor Medida-F (0,8826), e Sira e Dermason respondem pelas maiores confusões: 247 e 177 exemplos nos dois sentidos, seguidas por Barbunya e Cali (57 e 47).

\textbf{Comparação com o modelo linear.} A regressão logística chega a uma Medida-F de 0,9365 e perda de 0,2190. O ganho pareado da MLP é de $0{,}0070 \pm 0{,}0050$ na Medida-F, com vantagem em 8 das 10 partições. A MLP é melhor, com margem pequena, o que sugere que as fronteiras entre as classes são em boa parte quase lineares nesses atributos.

\subsection*{d. Conclusão da etapa}

A rede final generaliza de forma consistente, com Medida-F em torno de 0,94 e pouca variação entre partições. A estimativa é levemente otimista: como o enunciado prevê, a configuração foi escolhida com dados que também entram no k-fold, inclusive o teste da Questão 6.

Duas melhorias parecem mais promissoras do que aumentar a rede: atributos ou amostras adicionais voltados à separação entre Sira e Dermason, e parada antecipada em conjunto com calibração das probabilidades, já que a perda se mostrou mais sensível que a Medida-F. São hipóteses que este estudo não testou.

\FloatBarrier

\section{Conclusão geral}

A MLP aprendeu o problema com facilidade e chegou a uma Medida-F macro de $0{,}9436 \pm 0{,}0085$ em validação cruzada. Os principais achados de cada etapa foram:

\begin{itemize}
  \item A inicialização dos pesos afetou pouco o resultado final, e a ReLU foi tão boa quanto a tanh e muito mais rápida que a logistic.
  \item A taxa de aprendizado e o momento definem a velocidade por meio da taxa efetiva $\eta/(1-m)$. Valores altos trocam velocidade por overfitting rápido e, no extremo, levam ao colapso da rede.
  \item A topologia teve efeito pequeno. A variação entre as 30 redes é menor que a variação entre partições dos dados.
  \item Mais dados continuaram reduzindo a perda em toda a faixa testada, com retorno decrescente em Medida-F a partir de cerca de 40\% do treino.
  \item Dez atributos de forma substituem os 16 originais sem perda mensurável. Colunas de ruído causaram overfitting, e a importância por permutação engana quando os atributos são correlacionados.
  \item O critério de perda $\le 0{,}001$ só foi atingido por uma rede grande com LR alta, e com forte overfitting. Nesta base, em que as classes se sobrepõem, ele mede a memorização do treino.
\end{itemize}

A regressão logística ficou a 0,0070 da MLP, e o erro restante se concentra entre Sira e Dermason. Os dados indicam, portanto, que o ganho adicional deve vir mais da informação de entrada do que de redes maiores.

\subsection*{Limitações}

\begin{itemize}
  \item Três seeds por configuração nas Questões 2, 3 e 5 detectam diferenças grandes, mas não resolvem empates pequenos. Por isso as seleções seguem regras fixadas de antemão, e várias escolhas são equivalentes na prática.
  \item As topologias têm camadas de mesmo tamanho, e LR e momento foram escolhidos em uma única topologia, sem reavaliação conjunta.
  \item Na Questão 4, a quantidade de dados varia junto com o número de atualizações dos pesos.
  \item A seleção final usa o teste, como pede o enunciado, e o k-fold usa a base inteira. Os resultados finais são, assim, levemente otimistas.
  \item Os tempos valem para uma única máquina e servem apenas para comparações relativas.
\end{itemize}


\begin{thebibliography}{9}
  \bibitem{koklu2020} M. Koklu e I. A. Ozkan. Multiclass classification of dry beans using computer vision and
    machine learning techniques. \emph{Computers and Electronics in Agriculture}, 174:105507, 2020.
  \bibitem{uci} Dry Bean Dataset. UCI Machine Learning Repository, 2020.
    \url{https://doi.org/10.24432/C50S4B}.
  \bibitem{sklearn} F. Pedregosa et al. Scikit-learn: Machine learning in Python.
    \emph{Journal of Machine Learning Research}, 12:2825--2830, 2011.
  \bibitem{geron} A. Géron. \emph{Hands-On Machine Learning with Scikit-Learn, Keras, and TensorFlow}.
    3.\ ed. O'Reilly, 2022.
\end{thebibliography}

\end{document}
